\documentclass[12pt, a4paper]{report}

\usepackage[a4paper,left=3cm,right=2cm,top=2.5cm,bottom=2.5cm]{geometry}
\usepackage[hidelinks]{hyperref}

\begin{document}

\chapter{Introduction}

% introduction stuff


\chapter{Literature Review}

\section{A Novel Method for Detecting Credit Card Fraud Problems}

Du et al.~\cite{du2024novel} proposed a novel machine learning framework for credit card fraud detection that addresses the severe class imbalance present in financial transaction datasets. In credit card transaction data, legitimate transactions vastly outnumber fraudulent transactions, making it difficult for conventional machine learning models to accurately identify fraudulent activity. The authors note that conventional oversampling techniques such as the Synthetic Minority Oversampling Technique (SMOTE) may generate unrealistic or overly generalized minority-class samples. To overcome these limitations, the authors proposed a combined approach named AE-XGB-SMOTE-CGAN.

The study used a publicly available credit card transaction dataset containing 284,807 transactions collected over two days in September 2013, of which only 492 were identified as fraudulent. The dataset contains anonymized features, including PCA-transformed variables, along with a class label indicating whether a transaction is fraudulent or legitimate. Fraudulent transactions are represented by the label 1, while legitimate transactions are represented by 0.

The proposed AE-XGB-SMOTE-CGAN framework combines four major components: an Autoencoder, Extreme Gradient Boosting (XGBoost), SMOTE, and a Conditional Generative Adversarial Network (CGAN). SMOTE is initially used to address the severe class imbalance by generating additional minority-class samples. The information obtained from these samples is then transferred to the CGAN to generate more diverse synthetic samples. The Autoencoder subsequently learns a lower-dimensional representation of the transaction data, and the resulting feature representation is provided to the XGBoost classifier. The final classification is based on the probability produced by the XGBoost model.

The authors evaluated the proposed method using Accuracy (ACC), True Positive Rate (TPR), True Negative Rate (TNR), and Matthews Correlation Coefficient (MCC). The proposed approach was compared with conventional machine learning techniques, including XGBoost, K-Nearest Neighbors (KNN), and Random Forest. The results demonstrated that AE-XGB-SMOTE-CGAN achieved approximately a 2\% improvement in accuracy compared with XGBoost and KNN. Furthermore, when the classification threshold was set to 0.35, the proposed method achieved an MCC approximately 8\% higher than XGBoost and 30\% higher than KNN.

The main contribution of the study is the integration of oversampling, generative modelling, representation learning, and gradient-boosted classification into a single fraud detection framework. By combining these techniques, the proposed method improves the handling of highly imbalanced transaction data while achieving better fraud detection performance than the conventional approaches considered in the study.

Overall, the paper demonstrates that combining multiple data mining and machine learning techniques can improve the detection of fraudulent credit card transactions, particularly in datasets where fraudulent transactions represent only a very small fraction of the total data.

\section{An AutoEncoder Enhanced Light Gradient Boosting Machine Method for Credit Card Fraud Detection}
Ding et al.~\cite{ding2024autoencoder} proposed a hybrid approach named AEELG (AutoEncoder Enhanced LightGBM) for credit card fraud detection. This method addresses the limitations of single machine learning models, which struggle to accurately and efficiently identify anomalous transactions as financial datasets grow in scale and dimensionality. The authors highlight that standard approaches suffer heavily from the extreme class imbalance inherently present in credit card data, where legitimate transactions vastly outnumber fraudulent ones.

The primary evaluation of this study utilized a credit card fraud dataset containing 31 features. Furthermore, to test the model's robustness in high-dimensional spaces, the researchers validated their approach on a secondary dataset consisting of 200 features.

The AEELG framework integrates unsupervised feature reconstruction with a robust supervised classifier. Initially, the Synthetic Minority Oversampling Technique (SMOTE) is applied to balance the raw training data. An unsupervised AutoEncoder is then used to automatically learn optimal feature representations and extract deep, latent patterns without relying on labeled data. These enhanced features are subsequently passed into a LightGBM (Gradient Boosting Decision Tree) algorithm for the final classification. During the experimental setup, the AutoEncoder was trained to reconstruct features across different thresholds, ultimately selecting a threshold of 2 based on Matthews Correlation Coefficient (MCC) scores.

The performance of the AEELG framework was tested against 22 other algorithms and evaluated across nine different sampling techniques, including oversampling, undersampling, and mixed sampling. To objectively evaluate the imbalanced data, the authors utilized seven evaluation metrics: Precision, Recall, Accuracy, F-measure, Area Under the Curve (AUC), MCC, and Balanced Classification Rate (BCR).

The results demonstrated that AEELG significantly outperformed the other models. On the 31-feature dataset, the proposed method achieved a recall of 94.85\%, representing an approximate 10\% improvement over compared models, alongside a BCR score of 97\%. The main advantage of this model is its ability to inherit the benefits of both core components. It eliminates the heavy reliance on annotated data for optimal feature extraction while maintaining the fast, accurate classification capabilities of LightGBM. It successfully strikes a balance by accurately detecting true anomalies while keeping false positives remarkably low.

Despite its strong performance across comprehensive metrics, the authors noted that AEELG's precision decreased relative to some single models. This was identified as a direct trade-off of prioritizing recall via SMOTE and AutoEncoder-enhanced feature reconstruction. To address this open problem, future work proposes exploring the combination of AEELG with generative adversarial networks (GANs) to improve the quality of generated and reconstructed features, thereby enhancing classification precision.

Overall, the paper demonstrates that integrating unsupervised feature reconstruction with a robust supervised classifier and SMOTE preprocessing is a highly effective and scalable way to handle severe class imbalance in financial fraud detection.


\chapter{Comparative Analysis}

% comparison table/discussion


\chapter{Research Gap}

% gaps identified from the papers


\chapter{Conclusion}

% conclusion

\bibliographystyle{IEEEtran}
\bibliography{references}



\end{document}
